% Generated draft: requires mathematical and source review.
\documentclass[11pt]{article}
\usepackage[margin=25mm]{geometry}
\usepackage{fontspec}
\setmainfont{Latin Modern Roman}
\usepackage{amsmath,amssymb,mathtools,unicode-math}
\setmathfont{Latin Modern Math}
\usepackage{enumitem,needspace,xurl,hyperref}
\setlength{\parindent}{0pt}
\setlength{\parskip}{5pt}
\setlength{\emergencystretch}{3em}
\newcommand{\N}{\mathbb{N}}
\newcommand{\Z}{\mathbb{Z}}
\newcommand{\Q}{\mathbb{Q}}
\newcommand{\R}{\mathbb{R}}
\newcommand{\C}{\mathbb{C}}
\newcommand{\F}{\mathbb{F}}
\newcommand{\A}{\mathbb{A}}
\newcommand{\PP}{\mathbb{P}}
\newcommand{\E}{\mathbb{E}}
\newcommand{\eps}{\varepsilon}
\newcommand{\dd}{\,\mathrm{d}}
\begin{document}
\section*{UnsolvedMath ID: 30005529}
\textbf{Source problem code:} catalogue entry
\subsection{Definitions and mathematical statement}
Let \(\mathbb H=\{x+iy\in\C:y>0\}\) carry the hyperbolic metric \(\dd s^2=(\dd x^2+\dd y^2)/y^2\), and let \(\Gamma=\operatorname{PSL}_2(\Z)\). The modular surface \(M=\Gamma\backslash\mathbb H\) is understood as a hyperbolic orbifold. Its geodesic flow acts on its oriented unit tangent bundle, identified with \(X=\Gamma\backslash\operatorname{PSL}_2(\R)\), rather than on \(M\) alone. In the standard right-action convention, write
\[
 g_t(\Gamma u)=\Gamma u\begin{pmatrix}e^{t/2}&0\\0&e^{-t/2}\end{pmatrix},\qquad t\in\R.
\]
For comparison, one convention for the horocycle flow on the same space is
\[
 h_s(\Gamma u)=\Gamma u\begin{pmatrix}1&s\\0&1\end{pmatrix},\qquad s\in\R.
\]
These formulas fix notation; the supplied record does not specify the horocycle parametrization used in the intended comparison.

A usual framework for a symbolic coding of \(g_t\) is a shift \(\sigma\) on a space of symbolic sequences \(\Sigma\), a positive roof function \(r\), and a coding map from the suspension
\[
 \Sigma^r=\{(\omega,v):\omega\in\Sigma,\ 0\leq v\leq r(\omega)\}/\bigl((\omega,r(\omega))\sim(\sigma\omega,0)\bigr)
\]
to the relevant part of \(X\), intertwining suspension time translation with \(g_t\). The source does not prescribe the alphabet, admissibility conditions, regularity, coverage, multiplicity, or measure-theoretic properties of the coding; this framework explains the terminology without imposing those unspecified requirements.

\emph{Requested construction, as formulated in the supplied registry.} Construct a symbolic coding of the geodesic flow on the modular surface that recovers the gap statistics obtained from the corresponding horocycle-flow coding.

\emph{Unresolved qualification.} The supplied original excerpt asks for a geodesic-flow coding that ``gives the same,'' but omits the preceding comparison. Neither that excerpt nor the reconstructed statement specifies the horocycle coding, the objects whose gaps are measured, their ordering, their normalization, the limiting procedure, the reference measure, or the quantifiers governing the comparison. Consequently, the requested equality of gap statistics cannot be expressed as a fully determined mathematical condition from this record. The construction request above must retain this unresolved dependence on the missing source context; it does not assert equality of any particular distributions. The original excerpt motivates the request by suggesting that a coding might combine approaches discussed in talks by people referred to only as Jens and Matthew.
\subsection{Short English statement}
Seek a symbolic description of the modular geodesic flow that reproduces the gap statistics of an intended horocycle coding. The relevant coding and gap statistics are unspecified in the supplied record, so the comparison requires source clarification. The supplied open status and verification labels are curation claims; this is a statement-only draft requiring human mathematical review.
\subsection{Sources}
\begin{itemize}
\item Supplied local UnsolvedMath registry record, permanent ID \(30005529\), title \emph{Compatible Geodesic and Horocycle Symbolic Codings}, exact source problem code \texttt{catalogue entry}. The record assigns the year 2023 and names no proposer.
\item \emph{Group Actions and Harmonic Analysis in Number Theory} (2023), Oberwolfach Reports, pp.\ 1195--1262, DOI \texttt{10.4171/owr/2023/22}. Supplied source URL and literature-source label ``Oberwolfach problem report'': \url{https://doi.org/10.4171/owr/2023/22}. The supplied background describes Oberwolfach Reports as documenting mathematical workshops and their open-problem discussions. No exact page or named speaker for this problem is supplied.
\item Supplied literature-source label ``classical modular coding'': \url{https://arxiv.org/abs/math/9707215}. No authors, title, or precise supporting result are provided in the record.
\item Literature assessment checked at \texttt{2026-08-22}: the supplied curation labels the problem open and states that classical symbolic codings of modular geodesic flow exist, while a coding reproducing the report's specified gap statistics has not been established. This is supplied dated curation, not independently verified literature assessment. The background explicitly calls it dated literature triage derived from the supplied OWR curation and requests independent verification of time-sensitive status and cited claims.
\item The registry labels the statement \texttt{corrected\_verified} and states that it was reconstructed and checked against the official Oberwolfach report. These are supplied verification claims; no independent source visit or verification is asserted here.
\end{itemize}
\end{document}
